\documentclass{article}
\usepackage[utf8]{inputenc}
\usepackage[T1]{fontenc}
\usepackage[french]{babel}
\usepackage{geometry}
\usepackage{listings}
\usepackage{xcolor}
\usepackage{hyperref}

\geometry{a4paper, margin=2.5cm}

\definecolor{codegreen}{rgb}{0,0.6,0}
\definecolor{codegray}{rgb}{0.5,0.5,0.5}
\definecolor{codepurple}{rgb}{0.58,0,0.82}
\definecolor{backcolour}{rgb}{0.95,0.95,0.95}

\lstdefinestyle{mystyle}{
    backgroundcolor=\color{backcolour},
    commentstyle=\color{codegreen},
    keywordstyle=\color{magenta},
    numberstyle=\tiny\color{codegray},
    stringstyle=\color{codepurple},
    basicstyle=\ttfamily\footnotesize,
    breakatwhitespace=false,
    breaklines=true,
    captionpos=b,
    keepspaces=true,
    numbers=left,
    numbersep=5pt,
    showspaces=false,
    showstringspaces=false,
    showtabs=false,
    tabsize=2
}

\lstset{style=mystyle}

\title{Alanya Admin --- Fonctionnalité Broadcast / Diffusion}
\author{Plan technique}
\date{\today}

\begin{document}

\maketitle

\section{Objectif}

Permettre à un administrateur Alanya d'envoyer des messages (texte, image, vidéo) ou de créer des statuts (stories) à tous les utilisateurs ou à un sous-ensemble ciblé. Les messages sont envoyés depuis un compte système « Alanya » avec badge certifié. Les statuts sont visibles dans la section statuts de l'app mobile.

\section{Architecture}

\subsection{Compte Alanya (système)}

Créer un user avec un numéro réservé :
\begin{itemize}
    \item \texttt{nom} = « Alanya »
    \item \texttt{pseudo} = « alanya »
    \item \texttt{alanyaPhone} = numéro réservé (ex: 0000)
    \item \texttt{type\_compte} = 0 (apparaît comme user normal côté mobile)
    \item \texttt{avatar\_url} = logo Alanya
\end{itemize}

Ce compte est le \texttt{senderID} de tous les broadcasts.

\subsection{Tables SQL}

\subsubsection{Table \texttt{broadcast}}
\begin{lstlisting}[language=SQL]
CREATE TABLE broadcast (
  id              INT AUTO_INCREMENT PRIMARY KEY,
  sender_id       INT NOT NULL,       -- FK users (compte Alanya)
  created_by      INT NOT NULL,       -- FK users (admin qui diffuse)
  content         TEXT,
  type            TINYINT DEFAULT 0,  -- 0=text, 1=image, 2=video
  media_url       VARCHAR(255),
  target_type     VARCHAR(20) NOT NULL, -- 'all','country','specific'
  target_criteria JSON,               -- {"idPays":5} ou {"userIds":[1,2,3]}
  recipient_count INT DEFAULT 0,
  sent_at         DATETIME NOT NULL,
  FOREIGN KEY (sender_id) REFERENCES users(alanyaID),
  FOREIGN KEY (created_by) REFERENCES users(alanyaID)
);
\end{lstlisting}

\subsubsection{Table \texttt{broadcast\_recipient}}
\begin{lstlisting}[language=SQL]
CREATE TABLE broadcast_recipient (
  id              INT AUTO_INCREMENT PRIMARY KEY,
  broadcast_id    INT NOT NULL,
  alanya_id       INT NOT NULL,
  conversation_id BIGINT,
  message_id      BIGINT,
  sent_at         DATETIME NOT NULL,
  FOREIGN KEY (broadcast_id) REFERENCES broadcast(id),
  FOREIGN KEY (alanya_id) REFERENCES users(alanyaID),
  FOREIGN KEY (conversation_id) REFERENCES conversation(conversID),
  FOREIGN KEY (message_id) REFERENCES message(msgID),
  UNIQUE (broadcast_id, alanya_id)
);
\end{lstlisting}

\subsection{Migration}
\begin{lstlisting}[language=bash]
migrations/018_broadcast_tables.sql
\end{lstlisting}

\section{Endpoints API}

\subsection{POST /api/admin/broadcasts --- Créer et envoyer}

\textbf{Middleware:} \texttt{adminAuth}

\textbf{Request body:}
\begin{lstlisting}[language=JSON]
{
  "content": "Message de test",
  "type": 0,
  "mediaUrl": null,
  "targetType": "all",
  "targetCriteria": {}
}
\end{lstlisting}

\textbf{Logique:}
\begin{enumerate}
    \item Vérifier admin auth
    \item Résoudre les destinataires selon \texttt{target\_type} :
    \begin{itemize}
        \item \texttt{all} : tous les users non-bannis
        \item \texttt{country} : users par pays (\texttt{targetCriteria.idPays})
        \item \texttt{specific} : users spécifiques (\texttt{targetCriteria.userIds})
    \end{itemize}
    \item Pour chaque destinataire :
    \begin{itemize}
        \item Créer ou réutiliser la conversation 1-1 entre « Alanya » et le user
        \item Insérer le message dans \texttt{message}
        \item Mettre à jour \texttt{conversation.lastMessage}
        \item Incrémenter \texttt{conv\_participants.unreadCount}
        \item Envoyer FCM push via \texttt{notifyNewMessage()}
    \end{itemize}
    \item Enregistrer dans \texttt{broadcast} + \texttt{broadcast\_recipient}
    \item Retourner le résumé
\end{enumerate}

\subsection{GET /api/admin/broadcasts --- Historique}

\textbf{Query params:} \texttt{page}, \texttt{limit}

\textbf{Response:}
\begin{lstlisting}[language=JSON]
{
  "items": [...],
  "total": 42,
  "page": 1,
  "limit": 20
}
\end{lstlisting}

\subsection{GET /api/admin/broadcasts/:id --- Détails}

Retourne le broadcast + liste des destinataires avec statut d'envoi.

\section{Logique de conversation}

\begin{lstlisting}[language=JavaScript]
// Trouver ou créer la conversation 1-1
let conversation = await pool.execute(
  `SELECT c.conversID FROM conversation c
   JOIN conv_participants p1 ON p1.conversID = c.conversID
   JOIN conv_participants p2 ON p2.conversID = c.conversID
   WHERE c.isGroup = 0
     AND p1.alanyaID = ? AND p2.alanyaID = ?
   LIMIT 1`,
  [alanyaSenderId, recipientId]
);

if (!conversation.length) {
  // Créer la conversation
  const [conv] = await pool.execute(
    `INSERT INTO conversation (isGroup, lastMessage, lastMessageAt, lastMessageSenderID, lastMessageType)
     VALUES (0, ?, NOW(), ?, ?)`,
    [preview, alanyaSenderId, type]
  );
  conversationId = conv.insertId;

  // Ajouter les participants
  await pool.execute(
    `INSERT INTO conv_participants (conversID, alanyaID) VALUES (?, ?), (?, ?)`,
    [conversationId, alanyaSenderId, conversationId, recipientId]
  );
} else {
  conversationId = conversation[0].conversID;
  // Mettre à jour lastMessage
  await pool.execute(
    `UPDATE conversation SET lastMessage=?, lastMessageAt=NOW(), lastMessageSenderID=?, lastMessageType=?
     WHERE conversID=?`,
    [preview, alanyaSenderId, type, conversationId]
  );
}

// Envoyer le message
const [msg] = await pool.execute(
  `INSERT INTO message (senderID, conversationID, content, type, status, sendAt, mediaUrl)
   VALUES (?, ?, ?, ?, 1, NOW(), ?)`,
  [alanyaSenderId, conversationId, content, type, mediaUrl]
);

// Incrémenter unreadCount
await pool.execute(
  `UPDATE conv_participants SET unreadCount = unreadCount + 1
   WHERE conversID = ? AND alanyaID != ?`,
  [conversationId, alanyaSenderId]
);
\end{lstlisting}

\section{FCM Push}

Utiliser \texttt{notifyNewMessage()} existant :
\begin{itemize}
    \item Type : \texttt{message}
    \item Title : « Alanya »
    \item Body : preview du message
    \item ConversationId : l'ID de la conversation
    \item CallerId : l'ID du compte Alanya
\end{itemize}

\section{Statuts (Stories) broadcast}

En plus des messages, l'admin peut créer des statuts visibles par tous les utilisateurs.

\subsection{Logique de création}

\begin{lstlisting}[language=JavaScript]
// Créer un statut broadcast
const expiresAt = new Date(Date.now() + 24 * 60 * 60 * 1000);

const [statut] = await pool.execute(
  `INSERT INTO statut (alanyaID, type, text, mediaUrl,
   backgroundColor, createdAt, expiredAt)
   VALUES (?, ?, ?, ?, ?, NOW(), ?)`,
  [alanyaSenderId, type, content, mediaUrl,
   '#6366f1', expiresAt]
);

// Notifier tous les utilisateurs ciblés via FCM
const recipients = await resolveRecipients(targetType, targetCriteria);
for (const user of recipients) {
  await notifyNewMessage(user.fcmToken, {
    type: 'status',
    title: 'Alanya',
    body: content,
    statusId: statut.insertId,
  });
}
\end{lstlisting}

\subsection{Visibilité}

Par défaut, les statuts ne sont visibles que par les contacts mutuels (preferred contacts). Pour un broadcast statut :
\begin{itemize}
    \item « Tous » : insérer dans \texttt{statut\_views} pour TOUS les users non-bannis (ou modifier la logique côté mobile)
    \item « Par pays » : insérer pour les users du pays sélectionné
\end{itemize}

\textbf{Alternative plus propre :} Ajouter un champ \texttt{isBroadcast TINYINT DEFAULT 0} à la table \texttt{statut}. Côté mobile, si \texttt{isBroadcast = 1}, afficher le statut à tous sans vérifier les preferred contacts.

\subsection{Mobile}

Le statut broadcast apparaît dans la section statuts avec :
\begin{itemize}
    \item Avatar = logo Alanya
    \item Nom = « Alanya »
    \item Badge certifié
    \item Même UI qu'un statut normal, pas de différence visible
\end{itemize}

\section{Mobile --- UI spéciale}

Côté app Talky :
\begin{itemize}
    \item Détecter que le sender est le compte Alanya (\texttt{alanyaPhone} réservé)
    \item Afficher badge certifié (icône vérifiée)
    \item Photo de profil = logo Alanya
    \item Page profil : juste « Alanya » + icône certifiée, rien d'autre
    \item UI spéciale pour le message (fond différent, badge visible)
\end{itemize}

\section{Frontend Admin --- Composants}

\subsection{Page /broadcasts}
\begin{itemize}
    \item Historique des broadcasts envoyés
    \item Bouton « Nouveau broadcast »
    \item Colonnes : date, contenu (truncqué), cible, destinataires, statut
\end{itemize}

\subsection{Dialog Nouveau broadcast}
\begin{itemize}
    \item \textbf{Contenu} : textarea + upload image/vidéo
    \item \textbf{Ciblage} :
    \begin{itemize}
        \item Tous les utilisateurs
        \item Par pays (select dropdown)
        \item Utilisateurs spécifiques (recherche + sélection multi)
    \end{itemize}
    \item \textbf{Aperçu} : mini preview du message
    \item \textbf{Envoi} : bouton « Diffuser maintenant »
\end{itemize}

\subsection{Types TypeScript}
\begin{lstlisting}[language=TypeScript]
interface Broadcast {
  id: number;
  senderId: number;
  createdBy: number;
  content: string;
  type: number;         // 0=text, 1=image, 2=video
  mediaUrl: string | null;
  targetType: string;   // 'all' | 'country' | 'specific'
  targetCriteria: Record<string, unknown>;
  recipientCount: number;
  sentAt: string;
}

interface BroadcastRecipient {
  id: number;
  broadcastId: number;
  alanyaId: number;
  conversationId: number | null;
  messageId: number | null;
  sentAt: string;
}

interface BroadcastFormData {
  content: string;
  type: number;
  mediaUrl?: string;
  targetType: 'all' | 'country' | 'specific';
  targetCriteria: {
    idPays?: number;
    userIds?: number[];
  };
}
\end{lstlisting}

\section{Hooks React Query}
\begin{lstlisting}[language=TypeScript]
useBroadcasts(params)           // GET /admin/broadcasts
useBroadcastDetail(id)          // GET /admin/broadcasts/:id
useCreateBroadcast()            // POST /admin/broadcasts
\end{lstlisting}

\section{Fichiers à créer}
\begin{enumerate}
    \item \texttt{migrations/018\_broadcast\_tables.sql}
    \item \texttt{src/controllers/admin/broadcast.js}
    \item \texttt{src/routes/admin.js} (ajouter routes)
    \item \texttt{app/(dashboard)/broadcasts/page.tsx}
    \item \texttt{components/broadcasts/BroadcastDialog.tsx}
    \item \texttt{components/broadcasts/BroadcastHistory.tsx}
    \item \texttt{components/broadcasts/BroadcastPreview.tsx}
    \item \texttt{hooks/useBroadcasts.ts}
    \item \texttt{lib/mock-data.ts} (ajouter fonctions)
    \item \texttt{types/index.ts} (ajouter types)
    \item \texttt{app/(dashboard)/layout.tsx} (ajouter nav)
\end{enumerate}

\end{document}
